\section{Implementación y Ejecución de la Red Neuronal Convolucional}
La integración del modelo YOLOv2 se realiza mediante la implementación desarrollada por Milad 
Hasani \parencite{yolov2_implementation}. El script \textit{model.py} provee la funcionalidad 
para instanciar la arquitectura YOLOv2 Tiny y cargar sus pesos sinápticos, mientras que 
\textit{utils.py} proporciona funciones auxiliares relacionadas con la salida del modelo, 
tales como el cálculo de métricas y la generación gráfica de las cajas delimitadoras de 
predicción (\textit{bounding boxes}). Dichos scripts se ubicaron dentro del directorio 
\textit{TinyYoloV2} para desacoplar la implementación específica de la red del resto de 
los módulos (por defecto, Python trata las carpetas que contienen código como paquetes o 
módulos bajo el mismo nombre del directorio).
\medskip

El módulo de ejecución debe exponer la instancia del modelo, sus datos de entrada, la predicción 
resultante y un método para disparar la inferencia. La finalidad de esta estructura es permitir 
la configuración de los ganchos de ejecución posterior (\textit{forward hooks}) y la extracción 
de los metadatos de configuración. Este módulo se diseñó de la forma más independiente posible 
con el propósito de ser no intrusivo y adaptable a diversos entornos de ejecución. En otras
palabras, actúa como una abstracción de la implementación de un modelo.
\medskip

Para lograr esta abstracción, el módulo se estructuró mediante la clase \textit{ModelExecution}, 
la cual encapsula las estructuras de datos de la red y las funciones operativas de la 
implementación. Esta clase cuenta con un constructor que recibe como parámetro la ruta del archivo 
correspondiente a la imagen de entrada, realizando la carga y preprocesamiento de la imagen junto 
con los pesos del modelo pre-entrenado. Asimismo, el método \textit{execute} efectúa el paso hacia 
adelante (\textit{forward pass}) para obtener la inferencia y almacenar el resultado en un 
atributo de la instancia. Mediante este esquema es posible inicializar la red y rastrear su 
flujo de ejecución para acoplar los \textit{forward hooks} que capturan los tensores intermedios.